\chapter{经典解破裂以后：弱解与色散}
\label{chap:math-pde}

线性方程中的点源迫使我们把导数理解为分布。非线性方程即使从光滑初值出发，也可能自己产生跳跃。此时改变的不是守恒律，而是“解”的含义；方程仍应记录每个区域内的守恒以及穿越跳跃面的通量。

\section{从经典守恒律推弱形式}

设 \(u_t+f(u)_x=0\) 在 \(\mathbb R\times(0,\infty)\) 上有光滑解，初值为 \(u_0\)。取 \(\varphi\in C_c^\infty(\mathbb R\times[0,\infty))\)，乘方程积分。时间分部积分给
\[
\int_0^\infty\!\int_{\mathbb R}u_t\varphi\,\dd x\dd t
=-\int_0^\infty\!\int_{\mathbb R}u\varphi_t\,\dd x\dd t
-\int_{\mathbb R}u_0(x)\varphi(x,0)\,\dd x.
\]
时间无穷远的项因测试函数紧支而为零；空间分部积分也因其紧支没有边界项。把两项相加为零并移项，得到可在非光滑函数上继续使用的恒等式。

\begin{definition}[守恒律的弱解]
若 \(u,f(u)\in L^1_{\mathrm{loc}}(\mathbb R\times[0,\infty))\)，并且对每个上述测试函数都有
\[
\int_0^\infty\!\int_{\mathbb R}
\bigl(u\varphi_t+f(u)\varphi_x\bigr)\,\dd x\dd t
+\int_{\mathbb R}u_0(x)\varphi(x,0)\,\dd x=0,
\]
就称 \(u\) 为具有初值 \(u_0\) 的弱解。初值须以可使最后一项成立的迹意义进入。
\end{definition}

弱解定义保留了初值项，也写明了非线性通量的局部可积条件。若区域有物理边界，测试函数不在边界消失时还会出现边界通量；不能直接套用全空间公式。

\begin{theorem}[跳跃速度的 Rankine--Hugoniot 条件]
设 \(u\) 在直线 \(x=st\) 两侧分别为常数 \(u_L,u_R\)，通量 \(f\) 在这两值有定义。则 \(u_t+f(u)_x=0\) 在分布意义下成立，当且仅当
\[
s(u_R-u_L)=f(u_R)-f(u_L).
\]
\end{theorem}
\begin{proof}
把阶跃写成 \(u=u_L+(u_R-u_L)H(x-st)\)。上一章证明 \(H'=\delta_0\)；对复合变量求分布导数得到 \(u_t=-s(u_R-u_L)\delta(x-st)\)。同理 \(f(u)_x=(f(u_R)-f(u_L))\delta(x-st)\)。两者相加为零，恰等于给出的等式。反向也成立，因为跳跃曲线外两导数均为零。
\end{proof}

此条件只保证质量守恒，不保证弱解唯一。Burgers 方程的膨胀激波也满足这个等式，却与小黏性极限不符；要知道应留下哪一种弱解，还须考察熵不等式。

\section{Burgers 方程从光滑初值产生激波}

取无黏 Burgers 方程 \(u_t+(u^2/2)_x=0\)。在解尚光滑的区域，链式法则把它写成 \(u_t+uu_x=0\)。沿曲线 \(x(t)\) 若取 \(\dot x=u(x(t),t)\)，则
\[
\frac{\dd}{\dd t}u(x(t),t)=u_t+\dot x\,u_x=u_t+uu_x=0.
\]
所以一条特征线携带的 \(u\) 值不变。若初始点为 \(a\)，初值 \(u_0(a)\)，则 \(x=a+t u_0(a)\)、\(u=u_0(a)\)。这里 \(a\) 到 \(x\) 的映射必须一一对应，才能由位置反求唯一的 \(a\)。

取光滑有界初值 \(u_0(a)=-\tanh a\)。此时
\[
\frac{\partial x}{\partial a}=1+t u_0'(a)=1-t\operatorname{sech}^2a.
\]
当 \(t<1\) 时它处处为正，经典解单值；到 \(t=1\) 时在 \(a=0\) 变为零。隐式求导给 \(u_x=u_0'(a)/(1+t u_0'(a))\)，因而该处梯度趋于负无穷。初值一直光滑，失效的是光滑解的全局延续，并非初始资料有奇点。

\cref{fig:math-burgers-characteristics}把这一步画出来。每一条直线 \(x=a+t u_0(a)\) 是一条特征线，沿线 \(u\) 保持不变；特征线的疏密就是 \(u_x\) 的大小。\(\min u_0'=-\sech^2 0=-1\) 给出 \(t_*=1/(-\min u_0')=1\)，图中正是这条水平线附近，来自 \(a<0\) 与 \(a>0\) 的两族特征第一次相遇。在此之前，\(a\mapsto x\) 一一对应，可由位置反求初值点；此后同一位置对应多个初值点，单值的经典解不再存在。破裂的是解，不是初值。

\begin{figure}[htbp]
\centering
\includegraphics[width=.76\linewidth]{burgers_characteristics.pdf}
\caption{无黏 Burgers 方程取 $u_0(\xi)=-\tanh\xi$ 时的特征线。由于 $\min u_0'=-1$，理论预言 $t_*=1$；图中正是在这条水平线附近首先出现特征线交叉。}
\label{fig:math-burgers-characteristics}
\end{figure}

\section{弱解为何还要熵条件}

对左右常值 \(u_L,u_R\)，上一节的跳跃定理给 Burgers 激波速度 \(s=(u_L+u_R)/2\)。若 \(u_L>u_R\)，左侧特征速度 \(u_L\) 大于激波速度，右侧特征速度 \(u_R\) 小于它，两侧特征都进入激波。若 \(u_L<u_R\)，特征离开跳跃，出现“向外喷出信息”的膨胀激波；它满足弱形式，却不是稳定的小黏性极限。

\begin{definition}[凸熵不等式]
对守恒律 \(u_t+f(u)_x=0\)，若 \(\eta:\mathbb R\to\mathbb R\) 凸且可微，取满足 \(q'(u)=\eta'(u)f'(u)\) 的熵通量 \(q\)。弱解若对所选凸熵均满足 \(\eta(u)_t+q(u)_x\le0\) 的分布不等式，称熵解。不等式指它作用在每个非负紧支测试函数上非正。
\end{definition}

用最简单的膨胀跳跃算一次。取 \(u_L=0,u_R=1\)，则 \(s=1/2\)。对 Burgers 取 \(\eta(u)=u^2/2\)、\(q(u)=u^3/3\)，跳跃处的熵产生量为
\[
-s(\eta_R-\eta_L)+(q_R-q_L)
=-\frac12\cdot\frac12+\frac13=\frac1{12}>0.
\]
所以它违反熵不等式。相同初值的稀疏波解则在 \(0<x/t<1\) 取 \(u=x/t\)，左侧取零、右侧取一；把它代入光滑区方程并检查两条连接边界，可验证它是熵解。这个例子分清了两个问题：弱形式允许跳跃继续存在，熵条件选择哪一种继续方式。

小黏性为何挑出压缩激波，还可以直接计算。对 \(u_t+(u^2/2)_x=\nu u_{xx}\) 取行波 \(U(\xi)\)、\(\xi=x-st\)，要求 \(U(-\infty)=u_L>u_R=U(+\infty)\) 且两端导数趋零。把方程积分一次，并分别在两端代入，得到速度仍须是 \(s=(u_L+u_R)/2\)，而剖面满足
\[
\nu U'=\frac12(U-u_L)(U-u_R).
\]
记 \(\Delta=u_L-u_R>0\)，直接分离变量或代回可得
\[
U(\xi)=s-\frac\Delta2\tanh\!\left(\frac{\Delta\xi}{4\nu}\right).
\]
它单调下降，过渡层宽度与 \(\nu/\Delta\) 同阶；令 \(\nu\downarrow0\)，在跳跃线外逐点趋于 \(u_L,u_R\)。若试图让同一剖面从小值单调上升到大值，方程右侧在两端之间为负，左侧却要求 \(U'>0\)，矛盾。压缩与膨胀的区别因而来自一条可检验的常微分方程。

这个符号判断不依赖 Burgers 的特殊二次通量。设 \(f\in C^2\) 且 \(f''>0\)，即通量严格凸。对任意 \(u_L>u_R\)，Rankine--Hugoniot 速度是弦的斜率 \(s=[f]/[u]\)。凸性使 \(f'(u_R)<s<f'(u_L)\)：从左、右两侧来的特征都进入跳跃。更严格地，取任意两次可微凸熵 \(\eta\)，其通量 \(q\) 满足 \(q'=\eta'f'\)。令 \(a=u_R,b=u_L\)，并写
\[
 F(v)=f(v)-f(a)-s(v-a).
\]
函数 \(F\) 是凸函数与两端点连线之差，故 \(F(a)=F(b)=0\) 且 \(F(v)\le0\) 于 \([a,b]\)。跳跃处的熵产生量逐步化为
\[
 \begin{aligned}
 -s[\eta]+[q]
 &=\int_a^b\bigl(s-f'(v)\bigr)\eta'(v)\,\dd v
 =-\int_a^b F'(v)\eta'(v)\,\dd v\\
 &=\int_a^b F(v)\eta''(v)\,\dd v\le0.
 \end{aligned}
\]
最后一步的边界项因为 \(F(a)=F(b)=0\) 消失。若把左右值交换成 \(u_L<u_R\)，积分方向反转，严格凸熵使符号变正，因而膨胀跳跃被排除。这是从定义推出所有光滑凸熵的共同判据，不是只检查 \(\eta=u^2/2\) 的巧合。一般非凸通量不能照搬这一弦下方的图像；必须重新分析弦与曲线的位置。

被排除的膨胀跳跃由什么代替？对 \(a=u_L<b=u_R\)，寻找只依赖 \(\xi=x/t\) 的连续解 \(u(x,t)=U(\xi)\)。把它代入光滑区的守恒律，链式法则给 \((f'(U)-\xi)U'=0\)。在真的发生变化的区域，\(U'\ne0\)，所以 \(f'(U)=\xi\)。严格凸性使 \(f'\) 可逆，得到完整的稀疏波
\begin{equation}
 u(x,t)=
 \begin{cases}
 a,&x/t\le f'(a),\\
 (f')^{-1}(x/t),&f'(a)<x/t<f'(b),\\
 b,&x/t\ge f'(b).
 \end{cases}
 \label{eq:math-convex-rarefaction}
\end{equation}
在两条扇边，内部公式分别趋于 \(a,b\)，函数连续，故没有遗漏的跳跃 delta 项；扇内为经典解，所有光滑凸熵都满足等号。固定 \(x\ne0\) 令 \(t\downarrow0\)，它回到左右常值初态。对 Burgers，\(f'(u)=u\)，上式就在扇内给 \(u=x/t\)，与前面的具体计算一致。这里证明了两种显式解满足所述条件；一般初值下熵解的全局唯一性还需另证 \(L^1\) 收缩定理，不能由一次 Riemann 算例直接推出。

\begin{exercise}
对 \(u_L=2,u_R=0\) 求激波速度，计算 \(\eta=u^2/2\) 的跳跃熵产生量并判断符号；再把初值次序反过来说明为何同一跳跃公式不足以选解。
\end{exercise}


\section{黏性为何会挑出较平滑的解}

刚才给 Burgers 方程加上 \(\nu u_{xx}\)，并从行波中看见激波被摊成厚度约为 \(\nu/\Delta\) 的过渡层。要知道这个二阶项通常会做什么，先把非线性去掉，研究全空间的热方程 \(u_t=D u_{xx}\)，其中 \(D>0\) 的量纲是长度平方除以时间。采用前章的 Fourier 约定，\(\widehat{u_{xx}}(k)=-k^2\widehat u(k)\)，所以每个波数都满足常微分方程 \(\partial_t\widehat u=-Dk^2\widehat u\)，其精确解是
\[
 \widehat u(k,t)=\ee^{-Dk^2t}\widehat{u_0}(k).
\]
高波数 \(|k|\) 越大，衰减越快；这与后面色散项只改变相位而不衰减振幅的情形不同。对足够好的初值作逆变换，得到
\begin{equation}
 u(x,t)=\int_{\mathbb R}H_t(x-y)u_0(y)\,\dd y,
 \qquad H_t(x)=\frac{1}{\sqrt{4\pi Dt}}\exp\!\left(-\frac{x^2}{4Dt}\right).
 \label{eq:math-heat-kernel}
\end{equation}
这里 \(H_t\) 是点源初值的响应：它在任意正时间都是普通光滑函数，而 \(t\downarrow0\) 时在分布意义下趋于 \(\delta_0\)。这一极限可直接检查：\(\int H_t=1\)，代换 \(x=\sqrt{4Dt}\,z\) 后，对连续紧支测试函数 \(\varphi\) 有 \(\int H_t(x)\varphi(x)\,\dd x\to\varphi(0)\)。因此公式确实恢复初值，而不是只解了 \(t>0\) 的方程。

核非负且积分为一，所以若 \(u_0\ge0\)，则 \(u(t)\ge0\)；若 \(u_0\in L^1\)，交换积分还给 \(\int u(x,t)\,\dd x=\int u_0(x)\,\dd x\)。扩散保持总量，却摊开它的空间分布。对光滑且在无穷远衰减的解，用方程乘 \(\overline u\) 积分并分部积分，得到
\[
 \frac12\frac{\dd}{\dd t}\|u(t)\|_2^2
 =\operatorname{Re}\int\overline u D u_{xx}\,\dd x
 =-D\|u_x(t)\|_2^2\le0.
\]
若两个解有同一初值，把它们相减后重做这个估计，差的范数从零开始且不能增加，因此它们相同。公式 \(\ee^{-Dk^2t}\) 的模至多一，又使解算子能由光滑初值连续延拓到所有 \(L^2\) 初值。这样，Green 核、谱乘子、守恒和唯一性在同一个可算的 PDE 上相互检验；非线性方程的小黏性极限还需要额外估计，不能单凭热核的性质就宣布一般熵解唯一。

\section{色散怎样阻止波形无限变陡}

扩散把高频幅度衰减，色散则主要改变不同波数的相位。两者都能改写波形，但不能混为一谈。以 KdV 方程 \(u_t+6uu_x+u_{xxx}=0\) 为例，线性项的平面波 \(\ee^{\ii(kx-\omega t)}\) 满足 \(\omega=-k^3\)；不同 \(k\) 的相速度不同。若波振幅量级为 \(A\)、宽度量级为 \(L\)，非线性与三阶导数的量级分别为 \(A^2/L\) 和 \(A/L^3\)，平衡要求 \(A\sim L^{-2}\)。这只是尺度关系，尚未证明孤子存在。

为真正求解，令 \(u(x,t)=U(\xi)\)，\(\xi=x-ct\)，并要求 \(U,U',U''\to0\) 于 \(|\xi|\to\infty\)。代入原方程得 \(-cU'+6UU'+U'''=0\)。积分一次且用无穷远条件定积分常数为零，得到 \(U''=cU-3U^2\)。两边乘 \(2U'\) 再积分，第二个常数也由衰减条件确定为零：
\[
(U')^2=cU^2-2U^3=U^2(c-2U).
\]
对 \(c>0\)，函数 \(U(\xi)=\tfrac c2\operatorname{sech}^2(\sqrt c\,\xi/2)\) 满足此式，代回二阶方程也成立，故它是原 KdV 方程的行波解。振幅为 \(c/2\)，宽度与 \(c^{-1/2}\) 同阶，正好验证前面的尺度平衡。若只检验一次积分式而不回代，可能遗漏积分常数或错误分支。

KdV 也能从一个已写出误差阶的双向波模型中出现，而不必凭项的外形猜。设无量纲场满足
\[
u_{tt}-u_{xx}-\varepsilon\bigl[\alpha(u^2)_{xx}+\beta u_{xxxx}\bigr]
=O(\varepsilon^2),\qquad 0<\varepsilon\ll1,
\]
其中 \(\alpha,\beta\) 为固定实数。\(\varepsilon=0\) 时右行波依赖 \(x-t\)。若只观察它在长时间 \(t=O(\varepsilon^{-1})\) 的缓慢变化，令 \(\xi=x-t\)、\(\tau=\varepsilon t\)，并设 \(u=U_0(\xi,\tau)+\varepsilon U_1(\xi,\tau)+O(\varepsilon^2)\)。链式法则给 \(\partial_t=-\partial_\xi+\varepsilon\partial_\tau\)、\(\partial_x=\partial_\xi\)，故
\[
u_{tt}-u_{xx}=-2\varepsilon U_{0,\xi\tau}+O(\varepsilon^2),
\quad (u^2)_{xx}=(U_0^2)_{\xi\xi}+O(\varepsilon),
\quad u_{xxxx}=U_{0,\xi\xi\xi\xi}+O(\varepsilon).
\]
把这些式代入母方程的 \(O(\varepsilon)\) 系数，并在两端衰减的条件下对 \(\xi\) 积分一次，得到
\[
U_{0,\tau}+\alpha U_0U_{0,\xi}+\frac\beta2U_{0,\xi\xi\xi}=0.
\]
经变量和振幅的常数尺度变换，可化为上面的规范 KdV。这里的结论是对已给双向模型的形式渐近约化；若要保证近似在 \(t=O(\varepsilon^{-1})\) 上接近原解，还需对余项做稳定性估计，单靠逐项匹配不能推出严格误差界。

\section{窄带包络方程从哪一个小参数来}

另取一个明确的母方程 \(\ii\psi_t=\omega(-\ii\partial_x)\psi+g|\psi|^2\psi\)，其中实函数 \(\omega(k)\) 是线性色散关系，\(g\in\mathbb R\)。设波数集中在 \(k_0\) 附近宽度 \(O(\varepsilon)\)，振幅为 \(O(\varepsilon)\)，并写
\[
\psi(x,t)=\varepsilon A(X,T)\ee^{\ii(k_0x-\omega_0t)},
\quad X=\varepsilon(x-v_gt),\quad T=\varepsilon^2t,
\quad \omega_0=\omega(k_0),\quad v_g=\omega'(k_0).
\]
这里 \(\omega(-\ii\partial_x)\) 指 Fourier 空间中乘以 \(\omega(k)\) 的算子。若 \(\omega\) 在 \(k_0\) 附近三次可微，则在包络的窄带支撑内 Taylor 展开给
\[
\omega(k_0-\ii\varepsilon\partial_X)A
=\omega_0A-\ii\varepsilon v_gA_X
-\frac{\varepsilon^2}{2}\omega''(k_0)A_{XX}+O(\varepsilon^3).
\]
把表达式代回母方程，\(O(\varepsilon)\) 的载波项与 \(O(\varepsilon^2)\) 的群速度项分别抵消；在 \(O(\varepsilon^3)\) 得
\[
\ii A_T+\frac{\omega''(k_0)}2A_{XX}-g|A|^2A=0.
\]
这就是此母方程下的非线性 Schrödinger 包络方程。被舍弃的线性色散项在母方程残差中为 \(O(\varepsilon^4)\)，前提是包络的足够多导数有统一界；若演化使谱展宽到不再是 \(O(\varepsilon)\)，该估计就失效。KdV 的行波是精确解，包络方程则是带尺度条件的近似，两者的结论层级不同。

残差阶数也不是解的误差阶数。以只有线性自伴色散的方程为例，设近似解代回原方程留下 \(L^2\) 残差 \(r(t)\)，且真解与近似解初值相同。二者之差 \(e\) 满足 \(\ii e_t=\omega(-\ii\partial_x)e-r\)。自伴算子的演化保持 \(L^2\) 范数；把这个带源方程从零时刻积分，便有 \(\|e(t)\|_2\le\int_0^t\|r(s)\|_2\,\dd s\)。因此若 \(\|r\|_2\le C\varepsilon^4\) 且观察时间为 \(t\le T\varepsilon^{-2}\)，仅由积分就得到 \(\|e(t)\|_2\le CT\varepsilon^2\)，而不是 \(O(\varepsilon^4)\)。加入非线性后还要控制真解与近似解之差怎样反馈到非线性项；若没有这个稳定性估计，逐项匹配只能证明形式上的约化，不能声称长时间误差已受控。

\begin{exercise}
按上式逐项计算 \(\ii\partial_t\psi\) 和 \(\omega(-\ii\partial_x)\psi\) 到 \(O(\varepsilon^3)\)，指出群速度项为何消去；再将 KdV 行波解代回原方程直接检验。
\end{exercise}

